\chapter{Introducción}
\label{cap:introduccion}

En el mundo empresarial actual, la logística se ha consolidado como un factor determinante para el
éxito, la eficiencia operativa y la sostenibilidad de las organizaciones. La capacidad de optimizar
procesos, reducir costos, mantener la trazabilidad de la carga y garantizar la transparencia frente
a los clientes son exigencias fundamentales dentro de la cadena de suministro. Frente a estos
desafíos, la digitalización se ha convertido en una respuesta esencial: reemplaza los registros
manuales, centraliza la información y permite documentar las operaciones en el momento y el lugar
en que ocurren.

En este contexto surge \textbf{BeeTracer}, una solución tecnológica desarrollada por la empresa
chilena Bilix Ingeniería (Valparaíso) y orientada a la gestión, consulta y trazabilidad de faenas
logísticas, en particular las que se realizan en terminales extraportuarios. El sistema está
compuesto por una \textbf{aplicación móvil}, que permite al personal de terreno registrar datos y
evidencia fotográfica, y por una \textbf{plataforma web}, desde la cual se planifican, supervisan y
reportan las operaciones. Su propuesta de valor central es dejar un registro verificable del estado
de la carga en cada etapa crítica del proceso.

El presente informe corresponde a la segunda entrega de la asignatura Modelamiento de Software y
tiene por objetivo comprender el funcionamiento de BeeTracer mediante un proceso de
\textbf{ingeniería inversa}. No se dispone del código fuente ni de la documentación interna del
sistema, por lo que el modelo se construyó a partir de tres fuentes: la charla técnica dictada por
el equipo de BeeTracer \parencite{charlabeetracer2026}, la información pública de la plataforma
\parencite{bilix_sf} y la inferencia propia, aplicando las técnicas de la asignatura para deducir
la estructura que explica de forma coherente el comportamiento descrito.

Para ello se emplea el Lenguaje Unificado de Modelado (UML) \parencite{omg2017}, siguiendo el
enfoque de \textcite{larman2003}. A lo largo del informe se identifican los clientes, usuarios y
funciones del sistema; se representan sus funcionalidades mediante diagramas de casos de uso (de
alto nivel, expandidos y narrativos); se modela su comportamiento dinámico con diagramas de
secuencia y de colaboración; y se propone un diagrama de clases con su respectivo diccionario. Con
estas herramientas se busca ofrecer una visión clara, ilustrativa y técnica de las interacciones
que ocurren dentro del sistema.

Cabe precisar que las clases, atributos, operaciones y multiplicidades presentadas constituyen una
\textbf{propuesta de diseño inferida}, coherente con lo observado, pero no necesariamente idéntica
a la implementación real. Quedan fuera del alcance los módulos complementarios que la empresa
comercializa por separado (visor 3D de patios y aplicación de avisos de transporte) y el detalle de
la infraestructura física, por no contarse con información suficiente.
